% Options for packages loaded elsewhere
% Options for packages loaded elsewhere
\PassOptionsToPackage{unicode}{hyperref}
\PassOptionsToPackage{hyphens}{url}
\PassOptionsToPackage{dvipsnames,svgnames,x11names}{xcolor}
%
\documentclass[
  letterpaper,
  DIV=11,
  numbers=noendperiod]{scrartcl}
\usepackage{xcolor}
\usepackage{amsmath,amssymb}
\setcounter{secnumdepth}{-\maxdimen} % remove section numbering
\usepackage{iftex}
\ifPDFTeX
  \usepackage[T1]{fontenc}
  \usepackage[utf8]{inputenc}
  \usepackage{textcomp} % provide euro and other symbols
\else % if luatex or xetex
  \usepackage{unicode-math} % this also loads fontspec
  \defaultfontfeatures{Scale=MatchLowercase}
  \defaultfontfeatures[\rmfamily]{Ligatures=TeX,Scale=1}
\fi
\usepackage{lmodern}
\ifPDFTeX\else
  % xetex/luatex font selection
\fi
% Use upquote if available, for straight quotes in verbatim environments
\IfFileExists{upquote.sty}{\usepackage{upquote}}{}
\IfFileExists{microtype.sty}{% use microtype if available
  \usepackage[]{microtype}
  \UseMicrotypeSet[protrusion]{basicmath} % disable protrusion for tt fonts
}{}
\makeatletter
\@ifundefined{KOMAClassName}{% if non-KOMA class
  \IfFileExists{parskip.sty}{%
    \usepackage{parskip}
  }{% else
    \setlength{\parindent}{0pt}
    \setlength{\parskip}{6pt plus 2pt minus 1pt}}
}{% if KOMA class
  \KOMAoptions{parskip=half}}
\makeatother
% Make \paragraph and \subparagraph free-standing
\makeatletter
\ifx\paragraph\undefined\else
  \let\oldparagraph\paragraph
  \renewcommand{\paragraph}{
    \@ifstar
      \xxxParagraphStar
      \xxxParagraphNoStar
  }
  \newcommand{\xxxParagraphStar}[1]{\oldparagraph*{#1}\mbox{}}
  \newcommand{\xxxParagraphNoStar}[1]{\oldparagraph{#1}\mbox{}}
\fi
\ifx\subparagraph\undefined\else
  \let\oldsubparagraph\subparagraph
  \renewcommand{\subparagraph}{
    \@ifstar
      \xxxSubParagraphStar
      \xxxSubParagraphNoStar
  }
  \newcommand{\xxxSubParagraphStar}[1]{\oldsubparagraph*{#1}\mbox{}}
  \newcommand{\xxxSubParagraphNoStar}[1]{\oldsubparagraph{#1}\mbox{}}
\fi
\makeatother


\usepackage{longtable,booktabs,array}
\usepackage{calc} % for calculating minipage widths
% Correct order of tables after \paragraph or \subparagraph
\usepackage{etoolbox}
\makeatletter
\patchcmd\longtable{\par}{\if@noskipsec\mbox{}\fi\par}{}{}
\makeatother
% Allow footnotes in longtable head/foot
\IfFileExists{footnotehyper.sty}{\usepackage{footnotehyper}}{\usepackage{footnote}}
\makesavenoteenv{longtable}
\usepackage{graphicx}
\makeatletter
\newsavebox\pandoc@box
\newcommand*\pandocbounded[1]{% scales image to fit in text height/width
  \sbox\pandoc@box{#1}%
  \Gscale@div\@tempa{\textheight}{\dimexpr\ht\pandoc@box+\dp\pandoc@box\relax}%
  \Gscale@div\@tempb{\linewidth}{\wd\pandoc@box}%
  \ifdim\@tempb\p@<\@tempa\p@\let\@tempa\@tempb\fi% select the smaller of both
  \ifdim\@tempa\p@<\p@\scalebox{\@tempa}{\usebox\pandoc@box}%
  \else\usebox{\pandoc@box}%
  \fi%
}
% Set default figure placement to htbp
\def\fps@figure{htbp}
\makeatother


% definitions for citeproc citations
\NewDocumentCommand\citeproctext{}{}
\NewDocumentCommand\citeproc{mm}{%
  \begingroup\def\citeproctext{#2}\cite{#1}\endgroup}
\makeatletter
 % allow citations to break across lines
 \let\@cite@ofmt\@firstofone
 % avoid brackets around text for \cite:
 \def\@biblabel#1{}
 \def\@cite#1#2{{#1\if@tempswa , #2\fi}}
\makeatother
\newlength{\cslhangindent}
\setlength{\cslhangindent}{1.5em}
\newlength{\csllabelwidth}
\setlength{\csllabelwidth}{3em}
\newenvironment{CSLReferences}[2] % #1 hanging-indent, #2 entry-spacing
 {\begin{list}{}{%
  \setlength{\itemindent}{0pt}
  \setlength{\leftmargin}{0pt}
  \setlength{\parsep}{0pt}
  % turn on hanging indent if param 1 is 1
  \ifodd #1
   \setlength{\leftmargin}{\cslhangindent}
   \setlength{\itemindent}{-1\cslhangindent}
  \fi
  % set entry spacing
  \setlength{\itemsep}{#2\baselineskip}}}
 {\end{list}}
\usepackage{calc}
\newcommand{\CSLBlock}[1]{\hfill\break\parbox[t]{\linewidth}{\strut\ignorespaces#1\strut}}
\newcommand{\CSLLeftMargin}[1]{\parbox[t]{\csllabelwidth}{\strut#1\strut}}
\newcommand{\CSLRightInline}[1]{\parbox[t]{\linewidth - \csllabelwidth}{\strut#1\strut}}
\newcommand{\CSLIndent}[1]{\hspace{\cslhangindent}#1}



\setlength{\emergencystretch}{3em} % prevent overfull lines

\providecommand{\tightlist}{%
  \setlength{\itemsep}{0pt}\setlength{\parskip}{0pt}}



 


\KOMAoption{captions}{tableheading}
\usepackage{fancyvrb}
\usepackage{fvextra}
\RecustomVerbatimEnvironment{verbatim}{Verbatim}{breaklines=true,breakanywhere=true,fontsize=\small}
\makeatletter
\@ifpackageloaded{caption}{}{\usepackage{caption}}
\AtBeginDocument{%
\ifdefined\contentsname
  \renewcommand*\contentsname{Table of contents}
\else
  \newcommand\contentsname{Table of contents}
\fi
\ifdefined\listfigurename
  \renewcommand*\listfigurename{List of Figures}
\else
  \newcommand\listfigurename{List of Figures}
\fi
\ifdefined\listtablename
  \renewcommand*\listtablename{List of Tables}
\else
  \newcommand\listtablename{List of Tables}
\fi
\ifdefined\figurename
  \renewcommand*\figurename{Figure}
\else
  \newcommand\figurename{Figure}
\fi
\ifdefined\tablename
  \renewcommand*\tablename{Table}
\else
  \newcommand\tablename{Table}
\fi
}
\@ifpackageloaded{float}{}{\usepackage{float}}
\floatstyle{ruled}
\@ifundefined{c@chapter}{\newfloat{codelisting}{h}{lop}}{\newfloat{codelisting}{h}{lop}[chapter]}
\floatname{codelisting}{Listing}
\newcommand*\listoflistings{\listof{codelisting}{List of Listings}}
\makeatother
\makeatletter
\makeatother
\makeatletter
\@ifpackageloaded{caption}{}{\usepackage{caption}}
\@ifpackageloaded{subcaption}{}{\usepackage{subcaption}}
\makeatother
\usepackage{bookmark}
\IfFileExists{xurl.sty}{\usepackage{xurl}}{} % add URL line breaks if available
\urlstyle{same}
\hypersetup{
  pdftitle={annotations},
  colorlinks=true,
  linkcolor={blue},
  filecolor={Maroon},
  citecolor={Blue},
  urlcolor={Blue},
  pdfcreator={LaTeX via pandoc}}


\title{annotations}
\author{}
\date{}
\begin{document}
\maketitle


\paragraph{notes on BSI (2025)}\label{notes-on-bsi_wie_2025}

\begin{longtable}[]{@{}
  >{\raggedleft\arraybackslash}p{(\linewidth - 2\tabcolsep) * \real{0.0060}}
  >{\raggedright\arraybackslash}p{(\linewidth - 2\tabcolsep) * \real{0.9940}}@{}}
\caption{Table 1, BSI: annotations}\tabularnewline
\toprule\noalign{}
\begin{minipage}[b]{\linewidth}\raggedleft
id
\end{minipage} & \begin{minipage}[b]{\linewidth}\raggedright
annotations
\end{minipage} \\
\midrule\noalign{}
\endfirsthead
\toprule\noalign{}
\begin{minipage}[b]{\linewidth}\raggedleft
id
\end{minipage} & \begin{minipage}[b]{\linewidth}\raggedright
annotations
\end{minipage} \\
\midrule\noalign{}
\endhead
\bottomrule\noalign{}
\endlastfoot
1 & Wie KI unsere Sprache verändert -- Eine empirische Studie \\
2 & Die Grenze zwischen Werkzeug und Gegenüber, zwischen Medium und
Partner, wird dabei zunehmend durchlässig. Der Sprachgebrauch selbst
wird zum Ort der Aushandlung: zwischen individueller Intention und
algorithmischer Interpretation. Dieses neue Spannungsfeld erzeugt -- so
die zentrale These dieser Studie -- neue Sprachgewohnheiten, veränderte
semantische Ökonomien und möglicherweise sogar emergente Sprachformen,
die sich jenseits etablierter kultureller und linguistischer
Ordnungssysteme bewegen \\
3 & Es entsteht eine hybride Sprachpraxis, die Merkmale von
Alltagssprache, Programmiersprache, ökonomischer Instruktion,
metaphorischer Verdichtung und affektiver Andeutung kombiniert -- ein
Phänomen, das wir im weiteren Verlauf dieser Studie als Promptlinguistik
bezeichnen werden. \\
4 & Promptlinguistik \\
5 & Nutzer beginnen zunehmend, in einer Art Promptlogik zu denken und zu
sprechen. Diese Logik folgt nicht der klassischen Satzstruktur oder
narrativen Kohärenz, sondern ist geprägt von zielgerichteter
Instruktion, semantischer Kompression, elliptischen Formulierungen und
reduzierter Kontextualisierung \\
6 & Der Kommunikationsstil verschiebt sich von einem syntaktisch
vollständigen, sozialen Text hin zu einem effizienzorientierten, oft
kryptisch anmutenden Sprachgebrauch \\
7 & Diese Entwicklung lässt sich als Kognitionskompression begreifen --
ein psycholinguistischer Prozess, bei dem sprachliche Komplexität
reduziert wird, um mentale Entlastung, Output-Beschleunigung und
Interaktionsökonomie zu erreichen \\
8 & Prompten wird zur „sprechenden Skizze`` -- nicht mehr zur
Mitteilung, sondern zur Initialzündung algorithmischer Leistung. \\
9 & Prompting ist in vielen Fällen nicht mehr klar als Englisch, Deutsch
oder Spanisch identifizierbar, sondern als transkultureller,
maschinenadaptierter Idiolekt \\
10 & eröffnet zugleich neue Perspektiven auf eine mögliche
Globalgrammatik der Mensch-MaschineInteraktion. \\
11 & zielt die vorliegende Studie darauf ab, den Wandel sprachlicher und
kommunikativer Muster im Kontext der Nutzung generativer KI systematisch
zu erfassen und differenziert zu analysieren \\
12 & Nutzergruppe von 112 Heavy AI Usern, die über einen Zeitraum von
mindestens drei Monaten eine besonders intensive Interaktion mit
generativen Sprachmodellen --vorrangig ChatGPT -- aufwiesen \\
13 & Entsteht durch die Mensch-Maschine-Interaktion ein emergentes,
transkulturelles „Meta-Idiolekt``-- eine neue Ebene symbolischer Ordnung
zwischen Natur-, Kunst- und Programmiersprachen? \\
14 & Ziel ist es, ein erstes kognitiv-semiotisches Mapping der neuen
Sprachlandschaft im Zeitalter generativer KI zu entwerfen \\
15 & Menschen streben danach, bei gleichbleibender Ergebnisqualität
möglichst wenig kognitive Energie aufzuwenden (Kahneman, 2011). \\
16 & Erwartungshaltung, dass die Maschine semantische Lücken „mitdenkt``
, also kontextuelle Implikaturen und Nuancen selbstständig ergänzt \\
17 & kompensatorisches semantisches Vertrauen \\
18 & Diese Sprachveränderung steht nicht im Gegensatz zur Intelligenz
des Users, sondern ist Anpassung an das System: Sie basiert auf
gelerntem Promptingverhalten, das durch Erfahrung mit dem Modell
bestärkt wird. Die Sprache wird maschinen-kompatibel -- nicht mehr
menschlich-vollständig, sondern systemadäquat \\
19 & zeigen, dass bei intensiver Nutzung generativer KI die
durchschnittliche Promptlänge signifikant abnimmt, während die
semantische Dichte steigt \\
20 & kognitive Prozesse nicht auf das Gehirn beschränkt sind, sondern
sich über Werkzeuge und Umweltstrukturen erweitern können \\
21 & Generative KI kann unter dieser Perspektive als externes
semantisches Kognitionsorgan verstanden werden -- ein „sprachlicher
Cortex außerhalb des Körpers`` , der jedoch aktiv am Denkprozess
teilnimmt. \\
22 & Verarbeitungsflüssigkeit (Processing Fluency). \\
23 & kognitive Bevorzugung von flüssiger Information wirkt sich
nachweislich auf Urteile, Entscheidungen und Gedächtnisprozesse aus \\
24 & kein Zeichen sprachlicher Degradation, sondern Ausdruck einer
anpassungsfähigen Sprachintelligenz, die erkennt, dass komplexe Ästhetik
nicht belohnt wird --während funktionale Direktheit zu besseren
Ergebnissen führt. \\
25 & sprachliche Verlagerung vom Ausdruck zum Effekt -- das Ziel ist
nicht die Formulierung als solche, sondern die Antizipation eines
möglichst optimalen Outputs. Sprache wird zur strategischen Handlung,
nicht mehr zur sozialen Äußerung. \\
26 & Komplexe Idiome, Metaphern oder poetische Strukturen aus
spezifischen Sprachkulturen funktionieren weniger zuverlässig -- während
reduzierte, direktive, entkontextualisierte Sprache überdurchschnittlich
effektiv ist \\
27 & Symbolische Interaktionismus als soziologische Theorieperspektive
geht davon aus, dass soziale Wirklichkeit nicht objektiv vorliegt,
sondern in sozialen Interaktionen hergestellt, verhandelt und symbolisch
codiert wird \\
28 & Bedeutung nicht in Dingen selbst liegt, sondern im Prozess der
Bedeutungszuweisung durch Akteure. \\
29 & Sprache \\
30 & Hier deutet sich eine bislang unerforschte Form an: Symbolische
Interaktion ohne Intentionalität -- aber mit Wirkung. Der symbolische
Wert der Interaktion entsteht nicht aus dem Bewusstsein des Gegenübers,
sondern aus der Wiedererkennbarkeit, Kontinuität und Reproduzierbarkeit
interaktiver Schemata \\
31 & symbolischer Interaktionismus 2.0 \\
32 & In deutlicher Abgrenzung zur Vorstellung einer festgelegten
„Bedeutung der Wörter`` entwickelte Wittgenstein das Konzept der
Sprachspiele, das Sprache nicht als ein starres Abbild von Wirklichkeit,
sondern als regelgeleitete, kontextabhängige Praxis begreift \\
33 & Sprache ist demnach kein System aus Zeichen und Referenten, sondern
ein Handeln mit Zeichen, das sich an implizite Regeln knüpft, die
situativ ausgehandelt werden \\
34 & neue Form von Sprachspiel \\
35 & Ein Prompt ist kein Satz im klassischen Sinne, sondern ein
Sprechakt mit funktionaler Zielsetzung \\
36 & In der Interaktion mit generativer KI entwickelt sich zunehmend
eine neue Form davon: der Prompt-Idiolekt \\
37 & Die Bedeutung eines Prompt-Satzes erschließt sich nicht aus dem
Satz selbst, sondern aus der historischen Interaktionspraxis zwischen
User und Maschine. Damit entsteht eine Art privater Sprachspielraum, der
weder mit natürlicher Sprache noch mit Programmiersprache vollständig
deckungsgleich ist \\
38 & Die Bedeutung eines Prompts liegt nicht im linguistischen Gehalt,
sondern in seiner funktionalen Performanz im System \\
39 & Konzept des Cognitive Offloading \\
40 & Sprache wird zur Startfläche algorithmischer Leistung, nicht mehr
zur detaillierten Abbildung des eigenen Denkens. \\
41 & Syntaktisch komplexe Strukturen, wie sie etwa durch Nebensätze,
Konnektoren, explizite Subjekt-Prädikat-Relationen oder rhetorische
Absicherungen gekennzeichnet sind, weichen zunehmend minimalistischen,
kommandierenden oder referenzbasierten Sprachmustern \\
42 & inwiefern Nutzer das eigene Kommunikationsverhalten als funktional
transformiert wahrnehmen -- und ob sie bewusst oder unbewusst beginnen,
Sprache im Sinne einer kondensierten Steuerungslogik zu verwenden \\
43 & Es bildet sich ein idiosynkratischer Sprachraum aus, ein
Prompt-Idiolekt, der weder dem Standardregister der Alltagssprache noch
dem Code der maschinellen Sprache vollständig entspricht -- sondern
etwas Drittes ist: ein hybridisiertes Steuerungssystem, das durch
wechselseitige Eingewöhnung emergiert \\
44 & dialogischen Ko-Evolution mit einem System, das nicht intentional,
aber responsiv ist. Die KI lernt nicht individuell, aber sie reagiert in
berechenbarer Weise auf bestimmte semantische Muster -- und der Mensch
wiederum passt seine Sprache an diese Reaktionslogik an. Aus dieser
Rückkopplung entsteht eine Form funktional stabilisierter
Ausdrucksweisen, die zugleich persönlich markiert und maschinenadaptiert
ist \\
45 & Für Außenstehende sind diese sprachlichen Gebilde oftmals schwer
nachvollziehbar. Ihre Bedeutung erschließt sich nicht aus der Semantik
selbst, sondern aus der Interaktionserfahrung, die sie rahmt \\
46 & Gebrauchsmuster, die nur im konkreten Sprachspiel ihren Sinn
entfalten \\
47 & eine Art semantischer Gewöhnung, bei der sich Sprache zunehmend vom
sozialen Code der zwischenmenschlichen Kommunikation entfernt. \\
48 & Empirisch ließe sich diese Annahme über die Analyse wiederkehrender
Formulierungen, Prompt-Muster und sprachlicher Kürzel erfassen, ergänzt
durch Drittbeurteilungen hinsichtlich Verständlichkeit und
Kontextklarheit \\
49 & eröffnet damit den Zugang zu einem bislang kaum erforschten
Phänomen: der semantischen Selbstorganisation zwischen Mensch und
Maschine \\
50 & Processing Fluency \\
51 & das Erleben, Sprache nicht mehr im klassischen Sinne ``bilden'' zu
müssen, sondern sie lediglich als Steuerbefehl zu setzen. \\
52 & Interaktion mit generativer KI bei Nutzer zu einem gesteigerten
Erleben von Kommunikationsökonomie und subjektiver Sprachkontrolle
führt \\
53 & Gefühl manifestiert sich nicht nur in der Präferenz für knappe
Sprache, sondern auch in einer veränderten Erwartungshaltung gegenüber
Kommunikation generell: Sie soll funktionieren, nicht verweilen; sie
soll auslösen, nicht aushandeln; sie soll klar sein, nicht offen. \\
54 & psychologischen Zustand, in dem Sprache weniger als soziales Band
denn als symbolisches Steuerungsinstrument verstanden wird -- ein
Zustand, der nicht nur die KI-Nutzung selbst, sondern auch darüber
hinausgehende Kommunikationsgewohnheiten nachhaltig prägen kann. \\
55 & Ein deutschsprachiger Prompt wie „Mach es kürzer. Empathisch.
Bulletpoints. ``funktioniert in englischer, spanischer oder
französischer Übersetzung nahezu gleich gut. Die operative Sprachebene
ist nicht mehr idiomatisch, sondern algorithmisch organisiert. Damit
entsteht eine Form der Sprache, die sich über nationale Konventionen
hinwegsetzt -- eine emergente Maschinensprache, die sich durch ihre
technische Anschlussfähigkeit definiert, nicht durch ihren kulturellen
Resonanzraum. \\
56 & handelt es sich um eine funktionale, transkulturelle Sprache
zweiter Ordnung -- eine Sprache, die nicht dem zwischenmenschlichen
Verstehen, sondern der Optimierung maschineller Antwortlogik dient \\
57 & Ihre Grammatik ist eine Grammatik der Maschinenlesbarkeit, nicht
der sozialen Sensibilität \\
58 & plausibel anzunehmen, dass Nutzer:innen -- unabhängig von ihrer
jeweiligen Muttersprache -- ähnliche sprachliche Muster im Umgang mit
der KI entwickeln. \\
59 & potenzielle Entwicklung einer neuen, nicht-natursprachlichen, aber
humanadaptierten Sprachebene -- eine Sprache, die nicht Ausdruck
kultureller Differenz, sondern funktionale Antwort auf algorithmische
Konvergenz ist \\
60 & Hypothese 4: \\
61 & verweist auf einen weitreichenden Transfermechanismus: KI-basierte
Sprachgewohnheiten sind nicht an das Interface gebunden, sondern setzen
sich als kulturelle Technik durch \\
62 & zugrundeliegende Hypothese lautet deshalb, dass Menschen, die
generative KI intensiv nutzen, in anderen schriftsprachlichen Kontexten
eine deutlich erhöhte Neigung zu sprachlicher Verkürzung, funktionaler
Struktur und semantischer Direktheit zeigen \\
63 & Vielmehr handelt es sich um eine maschinenkompatible
Sprachökonomie, die nicht auf Selbstverwirklichung, sondern auf
funktionaler Output-Optimierung beruht \\
64 & Die Ausbildung individueller Prompt-Stile ist nicht lediglich ein
Nebenprodukt intensiver Nutzung, sondern ein zentraler Marker einer
neuen Mensch-Maschine-Kommunikation. \\
65 & Sprache wird nicht mehr in erster Linie als Werkzeug sozialer
Kohäsion begriffen, sondern als Mittel semantischer Beherrschung. Aus
dem Nutzer wird ein Prompt-Operator, aus dem Dialog eine
Steuerstruktur \\
66 & Dabei entsteht eine Art semantische „Lowest Common
Denominator``-Sprache: reduziert, anschlussfähig, funktional \\
67 & Bedeutung entsteht nicht mehr im sozialen Gebrauch einer bestimmten
Sprachkultur, sondern in der Rekurrenz funktionaler Semantik, die durch
maschinelles Lernen validiert wird. Die Maschine produziert keinen
kulturellen Sinn -- aber sie „belohnt``funktionale Sprache durch
qualitativ hochwertige Reaktionen. \\
68 & signifikante stilistische Verlagerung hin zu Kürze, Direktheit und
reduzierter Kontextualisierung -- selbst in Konstellationen, in denen
früher elaborierte oder höflichkeitsorientierte Sprache verwendet
wurde. \\
69 & Was im Maschinenraum der Sprache gelernt wurde, wirkt in den
sozialen Raum hinein -- durch Nachahmung, durch Habituation, durch
entlastende Wirksamkeit. \\
70 & Dabei wird Sprache nicht „verarmt`` , sondern umprogrammiert: weg
von kultureller Einschreibung, hin zu systemischer Wirkung. \\
71 & kündigt sich ein möglicher kultureller Paradigmenwechsel an: das
Verschwinden sprachlicher Höflichkeitsformeln, das Schrumpfen von
narrativen Einschüben, das Verschwinden emotionaler Zwischentöne \\
72 & Es zählt nicht mehr, wie etwas gesagt wird, sondern dass es
funktioniert. Der performative Charakter des Sozialen wird verdrängt
durch eine neue Grammatik der Funktionssprache \\
73 & Our findings should be taken as an invitation to revisit the common
assumption that LLMs respond best to lower perplexity prompts containing
simple and frequent words and grammatical structures. We find numerous
cases where LLMs do not learn from the data distribution in the way one
would assume based on perplexity scores and linguistic intuition \\
74 & Overall, our findings contradict Gonen et al.~(2022) and indicate
that the success of high or low perplexity prompts is particular to any
combination of dataset and model \\
75 & In general, we do not find that frequent synonyms lead to better
performance \\
\end{longtable}

\paragraph{notes on Yakura et al.
(2025)}\label{notes-on-yakura_empirical_2025}

\begin{longtable}[]{@{}
  >{\raggedleft\arraybackslash}p{(\linewidth - 2\tabcolsep) * \real{0.0083}}
  >{\raggedright\arraybackslash}p{(\linewidth - 2\tabcolsep) * \real{0.9917}}@{}}
\caption{Table 2, Yakura, LLM influence: annotations}\tabularnewline
\toprule\noalign{}
\begin{minipage}[b]{\linewidth}\raggedleft
id
\end{minipage} & \begin{minipage}[b]{\linewidth}\raggedright
annotations
\end{minipage} \\
\midrule\noalign{}
\endfirsthead
\toprule\noalign{}
\begin{minipage}[b]{\linewidth}\raggedleft
id
\end{minipage} & \begin{minipage}[b]{\linewidth}\raggedright
annotations
\end{minipage} \\
\midrule\noalign{}
\endhead
\bottomrule\noalign{}
\endlastfoot
76 & yakura-llm \\
77 & We detect a measurable and abrupt increase in the use of words
preferentially generated by ChatGPT---such as delve, comprehend, boast,
swift, and meticulous---after its release. These findings suggest a
scenario where machines, originally trained on human data and
subsequently exhibiting their own cultural traits, can, in turn,
measurably reshape human culture \\
78 & cultural feedback loop in which cultural traits circulate
bidirectionally between humans and machines \\
79 & the erosion of linguistic and cultural diversity, and the risks of
scalable manipulation \\
80 & the challenge is not if AI systems influence society, but how
profoundly and in which direction. \\
81 & The outcome is a linguistic and behavioral profile that, while
rooted in human language, exhibits systematic biases that distinguish it
from organic human communication \\
82 & ChatGPT exhibits a persistent preference for normative and socially
desirable communication patterns, such as politeness, neutrality, and
conflict avoidance---traits that are conventionally emphasized in
professional settings \\
83 & These tendencies reveal a self-consistent communicative style that
differentiates ChatGPT's language from spontaneous human interaction \\
84 & ChatGPT also demonstrates systematic lexical biases that reflect
its training data and optimization process \\
85 & these linguistic shifts stem primarily from direct machine-to-human
transmission---such as speakers reading from ChatGPT-co-authored
materials---rather than reflecting an internalized cultural trait \\
86 & We then found a clear post-ChatGPT increase in the usage of
delve---the word with the highest GPT score---rising in all categories
except Sports \\
87 & This suggests a two-phase diffusion process, where the words
preferred by the LLM first gains traction in fields with greater
exposure to content influenced by ChatGPT, such as in academic
environments, before filtering into more informal discourse \\
88 & These findings indicate that AI-driven linguistic shifts are not
confined to structured, scripted communication. Instead, they suggest a
plausible deeper process of linguistic adaptation, where certain words
become embedded in everyday speech rather than simply being imitated. \\
89 & This linguistic influence is not confined to domains where
LLM-generated text is integrated by early adopters---such as academia,
science, and technology---but it is spreading to other domains, such as
education and business. \\
90 & the mechanisms underlying this adoption remain unknown. This can be
the result of either direct imitation {[}51{]}, cognitive ease {[}52{]},
or deeper integration into human thinking processes {[}53{]}, or the
combination of all three. \\
91 & The uptake of words preferred by LLMs in real-time human-human
interactions suggests a deeper cognitive process at play---potentially
an internalization of AI-driven linguistic patterns. \\
92 & this measurable shift marks a precedent: machines trained on human
culture are now generating cultural traits that humans adopt,
effectively closing a cultural feedback loop. \\
93 & these traits no longer remain confined to interactions between
humans and AI systems but instead can diffuse further throughout
human-human communication. \\
\end{longtable}

\phantomsection\label{refs}
\begin{CSLReferences}{1}{0}
\bibitem[\citeproctext]{ref-bsi_wie_2025}
BSI, Brand Science Institute. 2025. {``Wie {KI} Unsere {Sprache}
Verändert -- {Eine} Empirische {Studie}.''}
\url{https://www.bsi.ag/cases/104-case-studie-wie-ki-unsere-sprache-veraendert---eine-empirische-studie.html}.

\bibitem[\citeproctext]{ref-yakura_empirical_2025}
Yakura, Hiromu, Ezequiel Lopez-Lopez, Levin Brinkmann, Ignacio Serna,
Prateek Gupta, Ivan Soraperra, and Iyad Rahwan. 2025. {``Empirical
Evidence of {Large} {Language} {Model}'s Influence on Human Spoken
Communication.''} arXiv.
\url{https://doi.org/10.48550/arXiv.2409.01754}.

\end{CSLReferences}




\end{document}
